\section{Overview}

Random matrix models are valuable tools for graph enumeration. As generating functions, they classify graphs by topology with weights set by the matrix size, $N$. This topological expansion is particularly relevant in the large $N$ limit in which planar graphs dominate. They have long been of interest in theoretical physics, where the sum over graphs can be interpreted as a discretised path integral for two-dimensional Euclidean quantum gravity. Through the use of multi-matrix models, one is able to endow graphs with statistical lattice models and hence study quantum gravity coupled to matter, or equivalently, bosonic string theory in a non-critical target space \cite{DiFrancesco:1993cyw, ginsparg1993lectures, Ambjorn:1997di}.

We consider a rather general class of matrix integrals of the form
\beq
Z = \int \rmd X_0 \rmd X_1 \ldots e^{-N \Tr V(X_0,X_1,\ldots)}
\eeq
where the $X_i$ are Hermitian $N\times N$ matrices with the $U(N)$ invariant Haar measure, and the potential is a polynomial function. One is typically interested in the moments of the model which can be captured by a generating function known as the resolvent. Here we define a \textit{solvable} theory to be one in which a closed expression for the resolvent can be determined \footnote{Note that this follows from the typical definition of solvability \cite{Kazakov:2000aq}}. The usual methods of solution can be broadly grouped into saddle-point methods, orthogonal polynomials, and the method of loop equations. There is, however, no straightforward mechanism for understanding whether a matrix model will be solvable and, for most models, direct analytic computation is virtually impossible.

With this in mind we revisit the Potts model on random graphs, first formulated by Kazakov \cite{Kazakov1988}. Since its inception this model has been studied a number of times as it furnishes a simple but non-trivial example of a multi-matrix integral that is not solvable by all the methods one typically brings to bear against one- and two-matrix models but the resolvent may still be determined. Recently Atkin, Niedner, and Wheater \cite{Atkin:2015ksy} were able to calculate resolvents corresponding to sums of random matrices using a variation on the saddle-point method. In particular, for the 3-state Potts model they were able to obtain results for not just the resolvent of a single matrix, but for the sum of two matrices and the sum of three matrices. This presents the question: can the loop equations method reproduce their findings?

This article is organised as follows: in Section 2 we establish the formalism we use to manipulate loop equations for general multi-matrix models. In Section 3 we describe an algorithmic approach for computing and systematically solving loop equations. We immediately apply this in Section 4 for the simple example of the two-matrix model with a general cubic potential. Following this we turn to a set of two-matrix models studied in \cite{} by the bootstrap method. We then turn to the 3-state Potts model and discuss the case of the fixed, mixed, and free boundary conditions. We provide a solution for the 3-state Potts model where the $S_3$ symmetry amongst spins is lifted establishing a result which was known but not written in the literature. 
